\documentclass[11pt]{article}
\usepackage{fontspec,amsmath,amssymb,amsthm,hyperref}
\usepackage[margin=1in]{geometry}
\setlength{\emergencystretch}{3em}
\title{Parameter KL by data processing and guarded real conversion}
\author{LangevinLean mathematical exposition}
\date{}
\begin{document}\maketitle
\noindent Source module: \texttt{IPLA/LSI/Marginal.lean}. Compile this proof document with XeLaTeX or LuaLaTeX. The accompanying HTML pairs the mathematical argument with the exact Lean statements and records the scope of the formalization.


The parameter reference law is exactly the projection of the actual lifted Gibbs target. Measurable projection decreases extended KL, and finite joint KL makes the real bound transfer valid. This measure-theoretic argument needs no marginal-density or strong-convexity formula.


\section{Define the projected target and prove its probability property}


With \(p=p_\theta\), define \(\Pi_N^\theta=p_\#\Pi_N\). The projection is a continuous linear map, hence measurable. For any probability measure \(\mu\), \((p_\#\mu)(\mathbb R^{d_\theta})=\mu(p^{-1}(\mathbb R^{d_\theta}))=\mu(\mathbb R^D)=1\). This proves probability preservation for any state law and then for the named target whenever \(\Pi_N\) is a probability. In the conditional LSI application that normalization comes from the LSI assumption itself.




\textbf{Correspondence with the displayed declarations.}


\begin{itemize}

\item \texttt{parameterTarget}: The parameter target is the actual lifted Gibbs law projected to its parameter coordinates. Probability normalization remains an explicit premise.
\item \texttt{parameterLaw\_isProbability}: Projecting a probability-valued state law preserves probability mass.
\item \texttt{parameterTarget\_isProbability}: The Gibbs parameter target is normalized whenever the actual joint target is normalized. No concentration or convexity hypothesis is introduced.
\end{itemize}


\noindent Lean declarations: \texttt{parameterTarget}, \texttt{parameterLaw\_isProbability}, \texttt{parameterTarget\_isProbability}.

\section{Extended KL decreases under the parameter projection}


The shared data-processing theorem for a measurable map yields

\[\operatorname{KL}(p_\#\mu\Vert p_\#\nu)\le\operatorname{KL}(\mu\Vert\nu).\]

At the density level this is convexity of \(u\log u\) under averaging over forgotten coordinates; the measure theorem does not require a constructed conditional density. If the right-hand divergence is finite, the inequality immediately excludes infinity on the left. Since both extended quantities are finite, order is preserved when taking their real values. This proves the real data-processing inequality with its necessary finiteness guard.




The parameter projection is a continuous linear map and hence measurable. Data processing gives the extended inequality \(\operatorname{KL}(p_\#\mu\Vert p_\#\nu)\le\operatorname{KL}(\mu\Vert\nu)\) before any real conversion. A finite joint KL therefore proves finite marginal KL, and only then can \texttt{ENNReal.toReal\_le\_toReal} transfer the order. This is a divergence statement about the actual pushed-forward target, with no density formula or optimizer concentration assumed.




\textbf{Correspondence with the displayed declarations.}


\begin{itemize}

\item \texttt{parameter\_kl\_le\_joint}: Extended-valued data processing for the common parameter projection.
\item \texttt{parameter\_kl\_ne\_top\_of\_joint}: Finite joint KL implies finite parameter KL.
\item \texttt{parameter\_klReal\_le\_joint}: The real data-processing inequality, with finiteness of both divergences established before taking their \texttt{toReal} values.
\end{itemize}


\noindent Lean declarations: \texttt{parameter\_kl\_le\_joint}, \texttt{parameter\_kl\_ne\_top\_of\_joint}, \texttt{parameter\_klReal\_le\_joint}.

\section{Transfer a full-joint or per-particle bound with the correct factor}


If finite full joint KL has real bound \(H\le B\), transitivity with real data processing gives parameter KL at most \(B\). If instead \(H/N\le B\) and \(N>0\), multiplication by the positive count yields \(H\le NB\), hence

\[\operatorname{KL}(p_\#\mu\Vert p_\#\nu)\le NB.\]

The proof uses the positive-denominator order equivalence before applying the full-joint transfer lemma. This factor is necessary because normalized joint entropy is not the same quantity as parameter-marginal entropy.




If the available estimate is intensive, \(H/N\le B\), positive \(N\) gives the full joint bound \(H\le NB\), and data processing then gives marginal KL at most \(NB\). It does not give marginal KL at most \(B\). The initial entropy in the final bound is still full joint entropy: projection decreases initial entropy, so replacing that term by initial parameter entropy would make the right-hand side smaller without a supporting argument.




\textbf{Correspondence with the displayed declarations.}


\begin{itemize}

\item \texttt{parameter\_kl\_bound\_of\_joint\_bound}: Any finite real bound on full joint KL is also a parameter-marginal bound.
\item \texttt{parameter\_kl\_bound\_of\_per\_particle\_bound}: A per-particle joint entropy bound controls unscaled parameter KL after multiplication by the particle count.
\end{itemize}


\noindent Lean declarations: \texttt{parameter\_kl\_bound\_of\_joint\_bound}, \texttt{parameter\_kl\_bound\_of\_per\_particle\_bound}.

\section{Specialize all transfers to the actual Gibbs parameter target}


Set \(\nu=\Pi_N\) in the three generic statements. By the definition of \(\Pi_N^\theta\), this yields extended data processing \(\operatorname{KL}(p_\#\mu\Vert\Pi_N^\theta)\le\operatorname{KL}(\mu\Vert\Pi_N)\), finiteness of parameter KL whenever joint KL is finite, and any real upper-bound transfer. The proofs are direct specializations and retain normalization of both laws and the finite-joint guard. No additional concentration or optimizer estimate appears in this measure projection step.




\textbf{Correspondence with the displayed declarations.}


\begin{itemize}

\item \texttt{parameter\_kl\_le\_gibbs\_joint}: Extended KL data processing with the actual Gibbs parameter target.
\item \texttt{parameter\_kl\_ne\_top\_of\_gibbs\_joint}: Finite actual Gibbs-relative joint KL gives finite parameter KL.
\item \texttt{parameter\_kl\_bound\_of\_gibbs\_joint\_bound}: A joint KL estimate transfers to the actual Gibbs parameter marginal, without constructing a formula for its marginal density.
\end{itemize}


\noindent Lean declarations: \texttt{parameter\_kl\_le\_gibbs\_joint}, \texttt{parameter\_kl\_ne\_top\_of\_gibbs\_joint}, \texttt{parameter\_kl\_bound\_of\_gibbs\_joint\_bound}.

\end{document}
